\chapter{ทบทวนความรู้พื้นฐานเรื่องจำนวนจริง}
\section{พหุนาม}
\subsection{การดำเนินการเกี่ยวกับพหุนาม}
\begin{conceptbox}[การบวกและการลบพหุนาม]
\textbf{การบวกและการลบพหุนาม}  
ให้บวกลบเฉพาะเอกนามที่เป็นชนิดเดียวกัน (มีตัวแปรและเลขชี้กำลังเท่ากัน)  
หากไม่สามารถบวกลบกันได้ ให้คงรูปเดิมไว้ เช่น
\begin{align*}
\left(2x^5+4x+5\right)+\left(x^5-x^2-2x\right)
&=3x^5-x^2+2x+5\\
\left(x^2-2x-1\right)-\left(2x^2-x-2\right)
&=-x^2-x+1\\
\left(3x^3-5x+2\right)+\left(-x^3+4x-7\right)
&=2x^3-x-5\\
\left(4x^2+x-6\right)-\left(x^2-3x+1\right)
&=3x^2+4x-7
\end{align*}
\end{conceptbox}

\begin{conceptbox}[การคูณพหุนาม]
\textbf{การคูณพหุนาม}  
ให้ใช้สมบัติการแจกแจง โดยคูณทุกพจน์ของพหุนามหนึ่งเข้ากับทุกพจน์ของอีกพหุนามหนึ่ง เช่น
\begin{align*}
6\left(x^2-x-1\right)
&=6x^2-6x-6\\
\left(x+3\right)\left(x^2-2x+1\right)
&=x^3+x^2-5x+3\\
\left(2x-1\right)\left(x^2+4x-3\right)
&=2x^3+7x^2-10x+3
\left(2x^2+4x+5\right)\left(x^2-2\right)\\
&=2x^4-4x^2+4x^3-8x+5x^2-10\\
&=2x^4+4x^3+x^2-8x-10\\
\end{align*}
\end{conceptbox}

\newpage

\begin{examplebox}[NETSAT คณิตศาสตร์ 2566]
กำหนด $A=-9 x^2+4 x+2$ และ $B=3 x^2+x-7$ ข้อใดต่อไปนี้ถูกต้อง
\begin{enumerate}
\item $2 A-B=-15 x^2-3 x-5$
\item $2 A-3 B=-21 x^2+24 x-17$
\item $A+3 B=7 x-19$
\item ไม่ถูกต้องทั้งข้อ 1, 2 และ 3
\end{enumerate}
\end{examplebox}

\newpage
\begin{examplebox}[NETSAT คณิตศาสตร์ 2566]
ข้อใดต่อไปนี้คือการกระจายนิพจน์ $x(x-8)-(x+2)(x-3)$ และทำให้อยู่ในรูปอย่างง่าย
\begin{enumerate}
\item $-7 x+6$
\item $-9 x+6$
\item $-10 x+6$
\item ถูกต้องทั้งข้อ 1, 2 และ 3
\end{enumerate}
\end{examplebox}

\newpage
\begin{examplebox}[NETSAT คณิตศาสตร์ 2566]
ข้อใดต่อไปนี้คือการแยกตัวประกอบของ $5 x+2 y-x y-10$
\begin{enumerate}
\item $(y-2)(5-x)$
\item $(x-2)(5-y)$
\item $(x-2)(y+5)$
\item $(y-5)(x+2)$
\end{enumerate}
\end{examplebox}

\newpage
\begin{remarkbox}[สูตรการกระจายที่ควรรู้]
\begin{itemize}
    \item[*] $(a+b)^2=a^2+2ab+b^2$
    \item[*] $(a-b)^2=a^2-2ab+b^2$
    \item[*] $(a+b)^3=a^3+3a^2b+3ab^2+b^3$
    \item[*] $(a-b)^3=a^3-3a^2b+3ab^2-b^3$
    \item[*] $a^2-b^2=(a+b)(a-b)$
    \item[*] $a^3+b^3=(a+b)(a^2-ab+b^2)$
    \item[*] $a^3-b^3=(a-b)(a^2+ab+b^2)$
\end{itemize}   
\end{remarkbox}

\begin{examplebox}[NETSAT คณิตศาสตร์ กุมภาพันธ์ 2565]
การกระจายนิพจน์ในข้อใดถูกต้อง
\begin{enumerate}
\item $(x-2)^2+x(x+4)=2 x^2+4$
\item $(x+4)^2-x(x-1)=7 x+16$
\item $4(x-1)^2+(x+2)^2=3 x^2-4 x$
\item ถูกต้องทั้งข้อ 1, 2 และ 3
\end{enumerate}
\end{examplebox}
\newpage
\begin{examplebox}[NETSAT คณิตศาสตร์ กุมภาพันธ์ 2565]
การแยกตัวประกอบของนิพจน์ในข้อใดถูกต้อง
\begin{enumerate}
\item $100 x^2-49 y^2=(10 x+7 y)(10 x-7 y)$
\item $27 x^3+125 y^3=(3 x+5 y)\left(9 x^2-15 x y+25 y^2\right)$
\item $x^3-12 x^2 y+48 x y^2-64 y^3=(x-4 y)^3$
\item ถูกต้องทั้งข้อ 1, 2 และ 3
\end{enumerate}
\end{examplebox}
\newpage
\begin{examplebox}[NETSAT คณิตศาสตร์ 2566]
ข้อใดต่อไปนี้คือการแยกตัวประกอบของนิพจน์ $25 x^2-81 y^2$
\begin{enumerate}
\item $(5 x-9 y)(5 x+9 y)$
\item $(5 x-81 y)(5 x+y)$
\item $(25 x-9 y)(x+9 y)$
\item ถูกต้องทั้งข้อ 1, 2 และ 3
\end{enumerate}
\end{examplebox}
\newpage
\begin{examplebox}[NETSAT คณิตศาสตร์ 2566]
ข้อใดต่อไปคือการกระจายพจน์ $(4 x-3 y)^3$ และทำให้อยู่ในรูปอย่างง่าย
\begin{enumerate}
\item $64 x^3-144 x^2 y+108 x y^2-27 y^3$
\item $64 x^3+144 x^2 y-108 x y^2-27 y^3$
\item $64 x^3-108 x^2 y+144 x y^2-27 y^3$
\item $64 x^3+108 x^2 y-144 x y^2+27 y^3$
\end{enumerate}
\end{examplebox}

\newpage
\begin{examplebox}[NETSAT คณิตศาสตร์ กุมภาพันธ์ 2567]
ข้อใดถูกต้อง
\begin{enumerate}
\item $(x+3)\left(x^2-3 x+9\right)=x^3-27$
\item $(x+3)(x+2)(x-2)(x-3)=x^4-25 x^2+36$
\item $(x+5)^2-(x+3)(x-3)=10 x+16$
\item $(x+3)(x-2)-x(x+1)=-6$
\end{enumerate}
\end{examplebox}

\newpage
\begin{examplebox}[NETSAT คณิตศาสตร์ กุมภาพันธ์ 2567]
ข้อใดต่อไปนี้แยกตัวประกอบได้ถูกต้อง
\begin{multicols}{2}
\begin{enumerate}
\item $a x^2-8 a x+16 a=a(x-4)^2$
\item $a^2+7 a+12=(a+4)(a+3)$
\item $a b-a-b+1=(a-1)(b-1)$
\item ถูกต้องทั้งข้อ 1, 2 และ 3
\end{enumerate}
\end{multicols}
\end{examplebox}
\newpage
\subsection{การหารสังเคราะห์}
\begin{examplebox}
จงใช้การหารสังเคราะห์เพื่อหาผลหารและเศษ จาการหารพหุนาม $3x^4+2x^2+3x-5$ ด้วย $x+1$
\end{examplebox}
\vspace{10cm}
\begin{examplebox}
จงใช้การหารสังเคราะห์เพื่อหาผลหารและเศษ จาการหารพหุนาม $x^3+x^2-18x+18$ ด้วย $x-3$
\end{examplebox}

\newpage
\begin{examplebox}
กำหนดให้ $p(x)=x^4-3x^3-5x^2+x-8$ และ $q(x)=x^2+4x-2$ จงใช้การหารสังเคราะห์เพื่อหาผลหารและเศษ จาการหารพหุนาม $p(x)$ ด้วย $q(x)$
\end{examplebox}

\newpage
\subsection{ทฤษฎีบทเศษเหลือ}
\begin{theorembox}
กำหนดให้ $p(x)$ เป็นพหุนาม และ $\alpha$ เป็นจำนวนจริง  
เมื่อหารพหุนาม $p(x)$ ด้วยพหุนาม $x-\alpha$ แล้ว  
\textbf{เศษจากการหารจะมีค่าเท่ากับ $p(\alpha)$}
\end{theorembox}

\begin{examplebox}
จงใช้ทฤษฎีบทเศษเหลือหาเศษเหลือจาการหารพหุนาม $x^4-5x^3+2x^2-x+2$ ด้วย $x-3$
\end{examplebox}
\vspace{3cm}
\begin{examplebox}
จงใช้ทฤษฎีบทเศษเหลือหาเศษเหลือจาการหารพหุนาม $x^3-x^2-x-15$ ด้วย $x-3$
\end{examplebox}
\vspace{3cm}
\begin{examplebox}
กำหนดพหุนาม $p(x)=x^3+6x^2-2\alpha x-3$ และ $q(x)=x^2+4x-2$ จงหาจำนวนจริง $\alpha$ ที่ทำให้ $p(x)$ ถูกหารด้วยพหุนาม $x+3$ ลงตัว
\end{examplebox}

\newpage
\begin{remarkbox}
กำหนดให้ $p(x)$ เป็นพหุนาม และ $\alpha$ เป็นจำนวนจริง
\begin{itemize}
    \item $p(\alpha)=0$ ก็ต่อเมื่อ พหุนาม $x-\alpha$ เป็นตัวประกอบของพหุนาม $p(x)$
    \item ถ้า $p(\alpha)=0$ จะเรียก $\alpha$ ว่า \textbf{ราก} หรือ \textbf{คำตอบ} ของพหุนาม $p(x)$  และของสมการ $p(x)=0$
\end{itemize}
\end{remarkbox}


\begin{examplebox}[NETSAT คณิตศาสตร์ กุมภาพันธ์ 2565]
ถ้าพหุนาม $x^3-x^2+a x+4$ ถูกหารลงตัวด้วย $x+2$ แล้ว $a$ มีค่าเท่าไร
\begin{multicols}{2}
\begin{enumerate}
\item $-2$
\item $2$
\item $-4$
\item $4$
\end{enumerate}
\end{multicols}
\end{examplebox}

\vspace{6cm}
\begin{examplebox}[NETSAT คณิตศาสตร์ กุมภาพันธ์ 2567]
ข้อใดคือเศษที่ได้จากการหารพหนุนามด้วย $3 x^3+4 x^2-6 x-10$ ด้วย $x-2$
\begin{multicols}{2}
\begin{enumerate}
\item $-6x^2$
\item $6 x$
\item $-6$
\item $6$
\end{enumerate}
\end{multicols}
\end{examplebox}

\newpage
\section{สมการตัวแปรเดียว $ax^2+bx+c=0$ เมื่อ $a\neq 0$}
\begin{theorembox}
สมการกำลังสองตัวแปรเดียว
\[
ax^2+bx+c=0 \quad (a\neq 0)
\]
จะมีคำตอบในระบบจำนวนจริงได้ไม่เกินสองคำตอบ  
โดยพิจารณาจากค่าของตัวแยกกำลังสอง (discriminant) คือ $b^2-4ac$ ดังนี้
\begin{itemize}
    \item ถ้า $b^2-4ac>0$ สมการจะมีคำตอบจริง \textbf{สองคำตอบที่แตกต่างกัน}
    \item ถ้า $b^2-4ac=0$ สมการจะมีคำตอบจริง \textbf{เพียงหนึ่งคำตอบ}
    \item ถ้า $b^2-4ac<0$ สมการจะ \textbf{ไม่มีคำตอบในระบบจำนวนจริง}
\end{itemize}
และคำตอบของสมการสามารถเขียนในรูป
\[
x=\dfrac{-b\pm\sqrt{b^2-4ac}}{2a}
\]
\end{theorembox}

\subsection{ผลบวกรากและผลคูณรากของสมการ $ax^2+bx+c=0$ เมื่อ $a\neq 0$}
\begin{remarkbox}
ถ้าสมการ $ax^2+bx+c=0$ มี $\alpha$ และ $\beta$ เป็นคำตอบ  
(อาจเป็นคำตอบซ้ำหรือคำตอบที่แตกต่างกัน) แล้วจะมีความสัมพันธ์ดังนี้
\begin{align*}
 \alpha+\beta &=-\dfrac{b}{a}\\
 \alpha\beta &=\dfrac{c}{a}
\end{align*}
\end{remarkbox}


\begin{examplebox}[NETSAT คณิตศาสตร์ กุมภาพันธ์ 2565]
ให้ $a$ และ $b$ เป็นคำตอบของสมการ $x^2+2 x-6=0$ ข้อใดถูกต้อง
\begin{multicols}{2}
\begin{enumerate}
\item $a+b=-2$
\item $(a-b)^2=28$
\item $a b=-6$
\item ถูกต้องทั้งข้อ 1, 2 และ 3
\end{enumerate}
\end{multicols}
\end{examplebox}

\newpage
\begin{examplebox}[NETSAT คณิตศาสตร์ สิงหาคม 2565]
ให้ $x=a$ และ $x=b$ เป็นคำตอบของสมการ $x^2+2 x-2=0$ ข้อใดต่อไปนี้ถูกต้อง
\begin{multicols}{2}
\begin{enumerate}
\item $a-b=2 \sqrt{3}$
\item $a+b=-2$
\item $a b=2$
\item ถูกต้องทั้งข้อ 1, 2 และ 3
\end{enumerate}
\end{multicols}
\end{examplebox}

\newpage
\begin{examplebox}[NETSAT คณิตศาสตร์ 2566]
ให้ $x=a$ และ $y=b$ เป็นคำตอบของของสมการ $(x+1)^2=7$ ข้อใดต่อไปนี้ถูกต้อง
\begin{multicols}{2}
\begin{enumerate}
\item $a=b$
\item $a+b>0$
\item $a b<0$
\item ถูกต้องทั้งข้อ 1, 2 และ 3
\end{enumerate}
\end{multicols}
\end{examplebox}

\newpage
\begin{examplebox}[NETSAT คณิตศาสตร์ 2566]
กำหนดให้ $x$ คือ คำตอบของสมการ $x^2-2 x-6=0$ ข้อใดต่อไปนี้ถูกต้อง
\begin{multicols}{2}
\begin{enumerate}
\item $x=1 \pm \sqrt{7}$
\item ผลรวมของคำตอบของสมการ คือ $2$
\item สมการนี้ไม่มีคำตอบ
\item ถูกต้องทั้งข้อ 1 และ 2
\end{enumerate}
\end{multicols}
\end{examplebox}